\subsection{Veroboard}
\markright{Veroboard}

\subsubsection*{Definition}
A veroboard (stripboard) is a semi-permanent prototyping board made of copper-clad
fiberglass (FR4) laminate \cite{wiki:breadboard}. Unlike a breadboard, component leads
are soldered to the copper tracks, producing a robust and reliable circuit suitable for
regular use. Horizontal copper strips run between rows of holes; where tracks must be
broken, a track-cutting tool removes a small section of copper.

\labimage{lab-01/assets/veroboard.jpg}%
  {Veroboard (stripboard) showing continuous copper tracks on the
   underside and the grid of plated-through holes on the component side.}{veroboard_lab}

\subsubsection*{Specifications}
Standard veroboard uses 2.54\,mm (0.1-inch) hole pitch, 1.6\,mm FR4 substrate, and
35\,$\mu$m (1\,oz) copper tracks. Copper strips carry several amperes continuously
\cite{analog:breadboard}.

\subsubsection*{Purpose of Use}
Veroboard bridges the gap between breadboard prototyping and PCB fabrication. It is
chosen when a circuit must be portable and reliable but producing a custom PCB is
unnecessary or uneconomical.

\subsubsection*{Applications}
Veroboard is used for audio amplifier assemblies, sensor interface boards, hobbyist
project builds, educational laboratory assignments requiring a permanent record, and
small-batch production of simple circuits \cite{boylestad2015}.
